% \chapter{Uncertainty Model for Systems of Twinned Systems (SoTS)}
% \label{chapter:uncertainty_model}

% Uncertainty in SoTS EventStreams
% 	Introduction
% 		Briefly recall the systematic study findings: few works directly address uncertainty.
% 		State purpose of this chapter: to classify, analyze, and formalize uncertainty in SoTS event streams, and show how it impacts downstream decision-making.
% 		Position: this chapter bridges state of the art (Ch. 3) → problem formulation (Ch. 5).
% 	Types of Uncertainty
% 		Aleatoric (inherent variability)
% 			random fluctuations in sensor readings (temp noise)
% 			Example - two DTs monitering the same beach water report slightly different readings due to calibration
% 		Epistemic Uncertainty (lack of knowledge)
% 			missing calibration, incomplete models, lack of training data
% 			Example - new DT node joins the SOTS with incomplete historical data
% 		Data Missingness	
% 			Little and Rubin taxonomy
% 				MCAR - missing at complete random (network packet loss)
% 				MAR - missing at random (sensor fails due to high load)
% 				MNAR - missing not at random (sensor consistently fails under high humidity)
% 			relate missingness to types of reliability risk in SoTS
% 		Figure: Uncertainty types (aleatoric, epistemic, missingness taxonomy).
% 	Hybird/Compounded Uncertainty
% 		SoTS uncertainty is rarely just one it, aleatoric, epistemic, and data missingness tend to interfact
% 		Example - intermittent connectivity + low fidelity DT leads to cascading uncertainty in CEP pipeline
% 	Sources of Uncertainty in SoTS streams - make into a table
% 		contituent level (DTs)
% 			sensor failures, calibration errors, model simplification
% 		communication level
% 			packet loss, latency, asynchrony between nodes
% 		sos level 
% 			dynamic membership - intermittent gaps
% 			interoperability mismatches - schema or sampling difference
% 		human/operational factors
% 			delayed updates, incorrect metadata, maintanance downtime
% 		Table: Sources of uncertainty → type → consequence.
% 	Propogation of Uncertainty
% 		show how uncertainty does not stay local but propogates upwards
% 			Sensor - DT model = mising sensor values distort DT fidelity
% 			DT - SOTS coordination = low fidelty DTs mislead collective state
% 			SOTS 0 CEO engineering = unreliable events streams reduce detection accuracy
% 		figure - pipeline diagram showing uncertainty flow (physical world - DT model - SoTS event stream - CEP patterns)
% 		cascading effect - one uncertain DT can degrade descisions at the system level 
% 		mention temporal compounding (uncertainty in earlier time steps degrades predictions downstream). That reinforces your justification for Kalman-based reconstruction.
% 	Why Uncertainty should be a first class concern
% 		In deterministic cep engines
% 			events are assumed as complete, accurate, and synchronous
% 			missing values are dropped
% 		in Dynamic SoTS
% 			this assumption breaks down continiously
% 			decisions without can be unsafe (triggering alarms, shutting down systems unnecessarily
% 		position CEP as critical descision making infrastructure, making uncertainty management not optional
% 	Data Quality frameworks and SoTS (REVAMP)
% 		Dimensions - accuracy, completeness, timeliness, consistency
% 		For SoTS constituents these map to 
% 			Directed - vulnerable to timelienss issues
% 			acknowledged - may drop out creating completeness gaps
% 			collaborative - rely on neighbours creating uncertainty in trust/consistency
% 		show that data quality frameworks are insuffient without uncertaint aware CEP
% 		Figure: Data quality framework mapped to DT/SoTS constituents.
% 	Consequences of SOTS level decision makinbg
% 		FP and FN in event detection
% 		loss of trust in system outputs
% 		inefficient resource allocation (mispredicted faults, unecessary maintanance
% 		example scenarios - smart city traffic control, smart energy grids
% 	CEP limitations under uncertainty	
% 		existing cep (esper, flink, drools fusionm)
% 			lack built in imputation, probabalistic reasoning or confidence annotations
% 			workaround - preprocess cep - brittle and not integrated
% 			need CEP that natively supports uncertainty as a part of the event semantics
% 	Summary
% 		SoTS are inherently uncertain
% 		uncertainty propogates and compounds affecting reliability
% 		current cep are not designed to handle this
% 		Transition: “The next chapter formulates the problem of designing an uncertainty-aware CEP framework and outlines the research questions guiding this work.”

        

% \section{Overview and Motivation}
% Uncertainty in Systems of Twinned Systems (SoTS) primarily arises when event streams are incomplete---that is, when expected sensor readings are missing, delayed, or inconsistent. 
% In large, distributed SoTS environments, constituent Digital Twins (DTs) depend on continuous event data to maintain accurate state synchronization and support reasoning through Complex Event Processing (CEP). 
% However, the autonomy of DTs, variability in network conditions, and transient sensor failures make \textit{missing data} an inevitable characteristic of SoTS event streams.

% This chapter focuses on \textbf{missing events as the primary form of uncertainty}, establishing a conceptual model that connects missingness to data reliability, reconstruction, and downstream reasoning. 
% It defines how missing data are represented, how uncertainty is expressed through confidence values, and how these confidence values evolve throughout the event pipeline---forming the basis for the uncertainty-aware architecture in Chapter~\ref{chapter:architecture}.

% % -------------------------------------------------------------
% \section{Nature of Missing-Event Uncertainty}
% \label{sec:nature_missing_uncertainty}

% Existing SoTS and Digital Twin frameworks emphasize reliability, synchronization, and interoperability 
% (e.g.,~\citep{rasheed2020digital, boschert2018digital, gabor2016concept}), 
% yet they commonly assume deterministic and complete data streams. 
% In practice, such completeness is rare: field-deployed IoT and cyber-physical systems---which underpin SoTS---frequently encounter missing or delayed measurements due to sensing or communication faults 
% (\citep{zhang2017missing, little2002statistical, ma2021handling}). 
% Consequently, the most direct and pervasive source of uncertainty in SoTS is \textbf{missing-event uncertainty}: the loss of expected data points within continuous streams.

% \subsection{Causes of Missing Events}
% \begin{itemize}
%     \item \textbf{Sensor-level interruptions.} Intermittent sensing due to calibration drift, energy-saving duty cycles, or environmental interference~\citep{shao2020missing, lin2020data}.
%     \item \textbf{Communication losses.} Network instability or congestion causing packet loss and transmission gaps~\citep{akdere2008comprehensive, stonebraker2005aurora}.
%     \item \textbf{System-level operations.} Constituent DTs may detach or suspend reporting for maintenance, a characteristic of autonomy in Systems-of-Systems~\citep{maier1998architecting, jamshidi2008systems}.
% \end{itemize}

% \subsection{Consequences}
% \begin{itemize}
%     \item \textbf{Temporal gaps and misalignment.} Missing readings violate the continuous cadence assumed by CEP engines~\citep{cugola2012processing, margara2014complex}.
%     \item \textbf{Propagation of incomplete context.} Data loss in one DT undermines dependent reasoning processes in others.
%     \item \textbf{Reduced reliability of event reasoning.} Traditional CEP frameworks ignore missingness, operating on deterministic event sequences~\citep{cugola2012processing}.
% \end{itemize}

% \begin{figure}[H]
% \centering
% \fbox{\includegraphics[width=0.8\textwidth]{figures/ch3_missing_sequence.pdf}}
% \caption{Illustrative sequence showing an observed stream with intermittent missing events.}
% \label{fig:ch3_missing_sequence}
% \end{figure}

% % -------------------------------------------------------------
% \section{Characterizing Missingness}
% \label{sec:missingness_characterization}

% To reason systematically about missing events, the framework adapts \citet{little2002statistical}'s taxonomy of missing-data mechanisms, interpreted for event streams (Table~\ref{tab:missingness_types}).

% \begin{table}[H]
% \centering
% \caption{Characterization of missingness types in SoTS event streams.}
% \label{tab:missingness_types}
% \begin{tabular}{p{2cm}p{8cm}p{4cm}}
% \toprule
% \textbf{Mechanism} & \textbf{Description in SoTS Context} & \textbf{Example} \\
% \midrule
% MCAR & Missing independent of data values or context. & Random packet drop. \\
% MAR & Missingness correlated with observable factors. & High humidity increases sensor failure. \\
% MNAR & Missing dependent on unobserved variables or system behavior. & Node intentionally suspends transmission. \\
% \bottomrule
% \end{tabular}
% \end{table}

% Each mechanism implies different recoverability and confidence potential when reconstructing events: MCAR data can be interpolated with minimal bias, while MNAR data introduce epistemic uncertainty that cannot be reduced without contextual reasoning.

% % -------------------------------------------------------------
% \section{Conceptual Representation of Uncertainty Through Confidence}
% \label{sec:representation_confidence}

% Rather than modeling full probability distributions for every event, this framework encodes uncertainty as a \textbf{confidence attribute} directly within each reconstructed event.

% \subsection{Event Augmentation Schema}
% Each event~$e$ is represented as:
% \[
% e = \langle id,\, t,\, \mathbf{x},\, c,\, s \rangle
% \]
% where:
% \begin{itemize}
%     \item $id$ --- event identifier,
%     \item $t$ --- timestamp,
%     \item $\mathbf{x}$ --- vector of observed or estimated attributes,
%     \item $c \in [0,1]$ --- confidence value,
%     \item $s \in \{\text{observed}, \text{reconstructed}\}$ --- source flag.
% \end{itemize}
% Observed events receive $c=1.0$; reconstructed events obtain $c<1.0$ based on reconstruction model reliability.

% \subsection{Confidence as Compressed Uncertainty}
% Full probabilistic tracking (e.g., covariance propagation) is computationally infeasible in real-time SoTS pipelines. 
% Confidence therefore serves as a compressed scalar summarizing:
% \begin{itemize}
%     \item local model prediction variance,
%     \item number of consecutive imputations, and
%     \item average residual error.
% \end{itemize}
% Confidence becomes an attribute accessible to CEP queries (e.g., \texttt{WHERE confidence > 0.7}), supporting uncertainty-aware reasoning.

% \begin{figure}[H]
% \centering
% \fbox{\includegraphics[width=0.8\textwidth]{figures/ch3_event_schema.pdf}}
% \caption{Augmented event structure with confidence attribute.}
% \label{fig:ch3_event_schema}
% \end{figure}

% % -------------------------------------------------------------
% \section{Reconstruction as Uncertainty Reduction}
% \label{sec:reconstruction_uncertainty}

% Reconstruction transforms data \textit{absence} into \textit{approximation}, thereby reducing uncertainty.

% \subsection{Role of Reconstruction}
% When the StreamManager detects a missing event, the Reconstructor estimates:
% \[
% \hat{x}_t = f(x_{t-1}, x_{t-2}, \dots)
% \]
% and derives a confidence~$c_t$ from model reliability (e.g., Kalman residual variance).

% \subsection{Reconstruction Methods}
% \begin{table}[H]
% \centering
% \caption{Common reconstruction methods and their confidence estimation bases.}
% \label{tab:reconstruction_methods}
% \begin{tabular}{p{3cm}p{6cm}p{5cm}}
% \toprule
% \textbf{Method} & \textbf{Mechanism} & \textbf{Confidence Basis} \\
% \midrule
% Kalman Filter & State-based prediction & Posterior variance or residual norm \\
% Linear Interpolation & Temporal linearity assumption & Fixed confidence decay with gap length \\
% Radial Basis Function & Localized interpolation & Error of nearest-neighbor fit \\
% \bottomrule
% \end{tabular}
% \end{table}

% \begin{figure}[H]
% \centering
% \fbox{\includegraphics[width=0.8\textwidth]{figures/ch3_reconstruction.pdf}}
% \caption{Illustration of missing-event reconstruction and corresponding confidence assignment.}
% \label{fig:ch3_reconstruction}
% \end{figure}

% % -------------------------------------------------------------
% \section{Uncertainty Propagation in Systems of Twinned Systems}
% \label{sec:uncertainty_propagation}

% Even after reconstruction, uncertainty can propagate through the hierarchical structure of a SoTS. 
% Each DT may consume data from others, allowing local uncertainties to influence higher-level reasoning and decision-making.

% \subsection{Hierarchical Propagation}
% SoTS architectures are typically multi-layered: DTs form the operational layer, and coordination engines form the supervisory layer~\citep{maier1998architecting, jamshidi2008systems}. 
% Uncertainty in constituent DTs therefore propagates upward through aggregation, similar to uncertainty spread observed in agent-based systems~\citep{macal2010tutorial, bonabeau2002agent}.

% \begin{figure}[H]
% \centering
% \fbox{\includegraphics[width=0.8\textwidth]{figures/ch3_propagation_layers.pdf}}
% \caption{Uncertainty propagation from sensor $\rightarrow$ DT $\rightarrow$ SoTS aggregate.}
% \label{fig:ch3_propagation_layers}
% \end{figure}

% \subsection{Mechanisms of Propagation}
% \begin{table}[H]
% \centering
% \caption{Mechanisms and analogues of uncertainty propagation in SoTS.}
% \label{tab:propagation_mechanisms}
% \begin{tabular}{p{3cm}p{7cm}p{3cm}}
% \toprule
% \textbf{Mechanism} & \textbf{Description} & \textbf{Analogue} \\
% \midrule
% Data aggregation & Combining multiple uncertain measurements. & Variance combination in inference. \\
% Model dependency & DT uses another DT's reconstructed stream. & Cascaded prediction error. \\
% Temporal chaining & Repeated reconstruction over time steps. & Error accumulation. \\
% Event composition & CEP rules depend on multiple uncertain events. & Probabilistic logic aggregation. \\
% \bottomrule
% \end{tabular}
% \end{table}

% For composite events $E_c$ derived from constituent events $e_i$ with confidences $c_i$:
% \[
% c(E_c) = 
% \begin{cases}
% \min(c_1,\dots,c_n) & \text{(conservative)}\\
% \frac{1}{n}\sum_i c_i & \text{(averaging)}
% \end{cases}
% \]

% \subsection{Implications for Event Reasoning}
% CEP operates on confidence-weighted event sets, analogous to probabilistic event reasoning~\citep{cugola2012processing, margara2014complex, cheng2018uncertain}. 
% Complex events with low constituent confidence can be deprioritized or flagged, enabling transparency in SoTS-level reasoning.

% \subsection{Relation to Agent-Based and Distributed Models}
% Agent-based models demonstrate how local uncertainty scales to emergent global behavior~\citep{macal2010tutorial, bonabeau2002agent}. 
% Similarly, each DT’s reconstruction confidence contributes to SoTS-level reliability. 
% If DT~$i$ provides state $x_i$ with confidence $c_i$, the system-level confidence for an aggregate metric $X$ can be expressed as:
% \[
% C_X = g(c_1, c_2, \dots, c_n)
% \]
% where $g$ may be a weighted mean or conservative minimum.

% \begin{table}[H]
% \centering
% \caption{Uncertainty propagation paths across SoTS hierarchy.}
% \label{tab:propagation_summary}
% \begin{tabular}{p{2cm}p{4cm}p{4cm}p{4cm}}
% \toprule
% \textbf{Level} & \textbf{Source} & \textbf{Propagation Path} & \textbf{Effect} \\
% \midrule
% Sensor & Missing data & Reconstructor & Local confidence $c_t < 1$ \\
% DT & Consumes reconstructed data & Model dependency & Inherited reduced confidence \\
% CEP & Combines events & Aggregation & Confidence attenuation \\
% SoTS & Aggregates DT outputs & Upward propagation & Emergent global uncertainty \\
% \bottomrule
% \end{tabular}
% \end{table}

% % -------------------------------------------------------------
% \section{Implications for Framework Design}
% \label{sec:design_implications}

% \begin{table}[H]
% \centering
% \caption{Mapping of uncertainty concepts to architectural mechanisms.}
% \label{tab:design_mapping}
% \begin{tabular}{p{5cm}p{8cm}}
% \toprule
% \textbf{Uncertainty Concept} & \textbf{Architectural Mechanism} \\
% \midrule
% Missing data introduces uncertainty & Reconstruction via imputation models. \\
% Need to preserve continuity & Tick-based scheduling of expected events. \\
% Uncertainty must be explicit & Confidence attribute added to event schema. \\
% Downstream consumers need flexibility & Separate topics for observed vs. reconstructed streams. \\
% Uncertainty affects reasoning & CEP integration with confidence-aware filters. \\
% Propagation across hierarchy & Confidence tracked from DT to SoTS-level reasoning. \\
% \bottomrule
% \end{tabular}
% \end{table}

% % -------------------------------------------------------------
% \section{Scope and Limitations}
% \label{sec:scope_limitations}
% This work focuses exclusively on uncertainty arising from missing data in SoTS event streams. 
% It assumes ordered event arrival and does not address late or out-of-order events, which would require event-time semantics and retraction capabilities at the CEP layer. 
% While these represent valuable future extensions, the present framework prioritizes continuous data reconstruction and confidence propagation as the primary mechanisms for uncertainty management.

% % -------------------------------------------------------------
% \section{Summary}
% \label{sec:ch3_summary}
% This chapter formalized missing events as the dominant source of uncertainty in Systems of Twinned Systems. 
% It introduced confidence as a lightweight event-level representation of uncertainty, detailed how reconstruction reduces and quantifies this uncertainty, and described how uncertainty propagates through the SoTS hierarchy. 
% These concepts directly inform the architecture presented in Chapter~\ref{chapter:architecture}, which implements the reconstruction, confidence propagation, and event dissemination mechanisms described here.
